\documentclass[10pt,a4paper]{article}
\usepackage[margin=23mm]{geometry}
\usepackage[T1]{fontenc}
\usepackage[utf8]{inputenc}
\usepackage{lmodern,amsmath,amssymb,booktabs,tabularx,microtype}
\usepackage[hidelinks]{hyperref}
\usepackage{xurl}
\setlength{\emergencystretch}{2em}
\renewcommand{\thesection}{A\arabic{section}}
\title{Online Resource 2: Supplementary Addendum\protect\\
\large The Review--Execution Gap in Agentic Software: Approval Binding and Recovery}
\author{Liang Hanghao, Xuan Wentao, Chen Qile, and Peng Peng}
\date{}
\begin{document}
\maketitle
\begin{center}\small
Automated Software Engineering manuscript\\
College of Computer Science and Electronic Engineering, Hunan University, Changsha, China\\
Corresponding author: Peng Peng, \texttt{hnu16pp@hnu.edu.cn}
\end{center}
\noindent\textbf{Scope.} This resource accompanies the manuscript on the review--execution gap. Online Resource 1 retains the original study documentation. The original 512-arm collection and separately frozen 117-pair deployment-A collection keep their original denominators, scores and inference. The analyses below are later diagnostics or separately frozen exploratory studies. Deterministic implementation checks are distinguished from hosted-agent observations. Sections A1--A2 report output-budget evidence, A3 examines receipt semantics, A4 reports local mechanism probes, A5 documents the historical audit, A6 maps controlled cases, A7 recounts preview exposure, and A8 reports configured ToolGate and public-reference workflows. Sections A9--A10 give the extended specification and design-selection guide; A11 documents native-approval and binding-scope reproduction; A12 reconstructs recovery availability and matched routes; A13 reports a separately frozen permission-state and feedback comparison; A14 tests native-to-store composition; A15 retains the detailed obligation-to-foundation map; A16 reports a separately frozen dynamic-recovery pilot; A17 reports a separately frozen check of reviewed-argument delivery across a process boundary in three pinned SDKs.

\section{Output truncation in original deployment B}
The diagnostic audit covered all 128 B tasks and 256 arms, counting each shared prefix once. It linked all 1,372 recorded provider responses to their events and saved inputs. Of these responses, 1,293 ended naturally with parseable actions and 79 reported \texttt{stop\_reason=max\_tokens}, with no parseable action. All 79 occurred after the shared review prefix: 76 had no text and three contained partial, unparseable text. The 4,096-token cap is established by the frozen configuration and adapter, not a retained HTTP wire capture.

\begin{table}[htbp]
\centering\small
\caption{Original B diagnostic counts. The original outcomes are retained.}
\begin{tabular}{lrr}
\toprule
Quantity & Context & Bound\\
\midrule
Arms & 128 & 128\\
Arms with at least one truncated suffix response & 33 & 28\\
Truncated suffix responses & 43 & 36\\
Original invalid-output terminations & 10 & 8\\
Original task completions & 117 & 119\\
\bottomrule
\end{tabular}
\end{table}

All eighteen invalid-output terminations ended in a truncated response after the history's one format-repair allowance had already been used. Five histories contained intervening parseable actions between their two truncated responses; these were not uniformly consecutive failures or transport retries. Among all 61 affected arms, 42 later completed, eighteen ended with invalid output, and one retained its missing-termination classification. No transport retry was recorded. Truncation at any point therefore cannot be equated with task failure.

The original bound-minus-context completion difference of two tasks coincides arithmetically with the unequal locations of these terminations: seven pairs have only a context termination, five only a bound termination, and three both. This is a descriptive decomposition, not an estimate of what the historical model would have done with a larger budget. The original failure scores, unknown handling and integrity violations remain intact; a later reporting failure does not undo a committed instance violation. The diagnostic supports describing these outcomes as truncation-associated invalid output, rather than naturally completed format noncompliance.

\section{Separately frozen output-budget sensitivity study}
Twenty-four tasks were sampled from the original task definitions, four from each of E, G, L, N, R and V. Each archived B review checkpoint and its visible prefix was continued under a contemporaneous $2\times2$ design: context/bound policy crossed with caps of 4,096/393,216 output tokens. The resulting 96 continuations were assigned before collection, with task blocks and Williams arm order fixed. There was one continuation per task and condition. This equal-scenario sample is not weighted like the original 128-task set.

Both budgets used the same current deployment, temperature zero, and extended waiting limits of 1,800 seconds per request and 21,600 seconds per continuation. The original twenty decision steps, twenty-four tool calls and one format-repair allowance per history were retained. The extended waiting limits also apply to the lower-cap condition; neither condition is an exact rerun of the original 90/720-second limits. The model alias does not establish that the immutable model version matched the historical collection.

\begin{table}[htbp]
\centering\small
\caption{Complete 96-continuation sensitivity results. $U$ and $I$ are scored separately; there are no unknown outcomes in this collection.}
\begin{tabular}{llrrrr}
\toprule
Policy & Output cap & $U=1$ & $I=1$ & Invalid output & Truncated/responses\\
\midrule
Context & 4,096   & 23/24 & 20/24 & 1 & 6/100\\
Bound   & 4,096   & 22/24 & 24/24 & 2 & 8/93\\
Context & 393,216 & 24/24 & 20/24 & 0 & 0/96\\
Bound   & 393,216 & 24/24 & 24/24 & 0 & 0/100\\
\bottomrule
\end{tabular}
\end{table}

The lower-cap conditions produced fourteen truncated responses across eleven arms. Three arms terminated with invalid output after truncation, before any tool action; they remain in the denominator and cost totals. In the expanded-cap conditions, all 196 responses ended naturally; thirteen reported more than 4,096 output tokens, with maximum 8,447. Thus, the larger cap was used in practice. All 389 requests received responses, with no outer transport retry.

The descriptive bound-minus-context completion difference is $-1/24$ at the lower cap and zero at the expanded cap. The difference of these differences is $+1/24$ (4.17 percentage points), supported by only three discordant tasks. No equivalence, non-inferiority or stable interaction claim follows. The four context integrity violations at each cap occurred at the programmed G handoff, not through naturally occurring agent substitution. Bound recovered with reviewed-candidate reuse in all four G tasks under each budget; those are two observations of the same four tasks, not eight independent tasks.

These observations show budget sensitivity in the sampled current continuations while retaining the integrity separation under both caps. They neither replace original results nor recover the counterfactual outcome of the eighteen historical failures. Actual continuation time and provider usage are available in the artifact report; low tool counts in prematurely terminated arms are not evidence of efficiency.

\section{Receipt unavailability versus uncertain commit outcome}
In the original R family, the agent receives a successful commit response before the first subsequent receipt query becomes unavailable. Authoritative state remains readable. A state-confirmation endpoint in that setting concerns an already acknowledged commit; it does not show that state inspection alone can attribute a commit whose success response was never delivered.

A separate deterministic check used the unchanged original Host, with durable copies of the full original request and key. Accept ($20\to20$) and Correction ($20\to21$), two policies and six history conditions produced 24 main records, plus eight controls. The fixed recovery script retried the original arguments and key once and then read full public state. It received only the client journal and public returns; injected history and private state were used only by the environment and scorer.

\begin{table}[htbp]
\centering\small
\caption{Deterministic recovery witnesses. Each row contains four records. These are not LLM-agent outcomes.}
\begin{tabularx}{\textwidth}{@{}Xll@{}}
\toprule
History & Correct own-request report & New recovery transitions\\
\midrule
Initial request not delivered & Confirmed after retry & One per record\\
Commit persisted, success reply lost & Historical success & Zero\\
Competing key commits same value & Rejected; goal true & Zero\\
Own commit, then reviewer disabled & Historical success & Zero\\
Own commit; retry delivery blocked & Unknown; goal true & Zero\\
Competing commit; retry delivery blocked & Unknown; goal true & Zero\\
\bottomrule
\end{tabularx}
\end{table}

There were twelve confirmed, four rejected and eight unknown reports, with one business transition in every final state. Four paired hidden histories had identical full client journals and public observations, including values, versions, sources and capabilities, but different request ownership. The competing request used the same candidate and grant with another key; it was not an independently approved actor. These pairs are concrete witnesses that a matching state need not identify the request that committed. They do not prove that all possible interfaces are permanently insufficient.

Controls also distinguished historical receipts from current permission: the original successful response remained replayable after reviewer disablement without a new transition, and a cached refusal remained replayable after reviewer restoration while a valid fresh request could succeed. Changed arguments under a previously used key were rejected as a collision; such a collision is not evidence that the original request failed.

Both runs passed all 316 implementation checks; rebuilding in another directory was not a new statistical sample. The result depends on a complete retained client request, retained idempotency records and complete SQLite state checkpoints. It does not cover partial persistence, process crashes during commit, cache eviction, lost client journals, real concurrent races or cross-service effects. Existing retry and receipt mechanisms suffice under these tested conditions. Agent ability to choose or interpret those mechanisms was not measured.

\section{Strong-binding overlap and local integration}
The two comparisons below preserve the native protections and avoid treating an unavailable abstraction as a failed system. They are local exploratory checks, not a complete external-system benchmark or an efficacy ranking.

\subsection*{Continuity Kernel model probe}
The unmodified author model at commit \nolinkurl{74841530988f1324bd927da15519486509372917} was called through its native submission interface. The matrix planned 24 records, with context, bound and the CK model for two value relations and four scenarios. Twenty-two records executed; two CK instance-substitution cells were retained as not representable because its exposed abstraction lacks a distinct human-reviewed candidate object. They were not counted as failures, accepted substitutions or zero integrity scores.

All eighteen records in the six jointly expressible cells agreed on coarse admission decisions for normal submission, changed evidence and changed predecessors. CK's Accept/Correction values were adapter encodings, and its evidence validity was a supplied predicate; these are not native demonstrations of the full correction or evidence-byte contract. Fourteen initial rejections had scripted legal recovery witnesses, all twenty-two executable requests replayed without another authoritative change, and eleven controls checked retained protections. The 150 assertions are implementation checks, not independent samples. This probe supports capability mapping, not superiority or absence of instance-binding capability in CK.

\subsection*{Native saved-plan workflow}
The application rule was an approval for a named logical saved-plan artifact, with a fresh approval required for another plan. HashiCorp's official automation guidance describes tying approval to the correct plan. OpenTofu treats a supplied saved plan as approved input, so both comparison routes added the same explicit trusted approval record, artifact identifier and digest checks before invoking the unmodified OpenTofu 1.13.1 binary. Those approval checks belong to the workflow wrapper, not a native human-approval API.

One route used saved-plan binding; the other retained the same binding and added explicit correction receipts. Both had the same original proposal, plan and approval information, and the baseline could reconstruct the relation from its logs. Accept and Correction crossed six scenarios and two routes to give 24 main records (twelve route pairs), plus eight controls. The local \texttt{terraform\_data} fixtures used no cloud resources, external providers, human reviewers or LLM calls.

Both routes satisfied the tested common obligations. Each route's twelve main records used 128 native CLI calls and sixteen approval records, including authorized state advancement in stale-state fixtures. Twelve scripted recoveries succeeded across both routes. All twenty-four main replays were refused as stale without another state change: this is stable non-reexecution, not a returned cached success receipt. The extension generated sixty extra receipt files totalling 40,521 bytes. The 240 checks and new-directory rebuild verify this constructed integration; they do not establish minimal storage cost, production performance or audit time saved.

These comparisons locate the contribution in the explicit four-role correction relation, successor and field-source obligations, and the original online evidence on agent continuation. Identity binding and recovery have strong existing foundations. The OpenTofu mapping did not port the full successor/provenance contract, and neither comparison demonstrated additional execution protection or audit efficiency over a strong binding baseline. A broader system comparison or an audit-task study would require a separately specified observable benefit.

\section{Historical correction-chain audit protocol}
The main text reports that an audit of 1,260 historical databases identified 84 retained machine-origin correction chains. This section records the population, discovery rule and per-chain counts so that the claim can be checked from the deposited evidence.

\textbf{Population and discovery.} The audit opened all 1,260 archived databases read-only. Each database belongs to one archived run of the restricted-agent benchmark (three requested configurations, 420 runs each). A correction chain was enumerated from the persisted transition table as a committed transition whose source kind is \texttt{AI\_SUGGESTION} and whose committed value differs from the proposal value stored on its candidate certificate. No task-success label was used for selection. All such transitions in the population were audited; none was excluded after inspection.

\begin{table}[htbp]
\centering\small
\caption{Historical correction-chain audit counts. Units are persisted chains and checks, not task successes.}
\begin{tabular}{lr}
\toprule
Quantity & Value \\
\midrule
Databases scanned & 1,260 \\
Correction chains found & 84 \\
Chains with an audited violation & 0 \\
Chains with no audited violation & 84 \\
Behavior-evaluable chains & 81 \\
Evidence associations checked & 252 \\
Candidate-certificate checks & 252 \\
Record versions checked & 168 \\
P6 provenance traces checked & 168 \\
Chains with unchanged-field checks & 0 \\
\bottomrule
\end{tabular}
\end{table}

Configuration counts are $G1=29$, $G2=26$ and $D1=29$; prompt-variant counts are $V1=27$, $V2=29$ and $V3=28$. Terminal classes of the containing runs are 80 \texttt{PASS\_UTILITY\_SAFE}, one \texttt{PASS\_SAFE\_BUT\_UTILITY\_FAILED}, two \texttt{MODEL\_API\_FAILURE} and one \texttt{TIMEOUT}. All 84 chains are single-field, so the audit supplies zero nonvacuous unchanged-field checks; copy-forward is supported by the controlled multi-field evidence, not by this audit.

\textbf{What the audit does and does not establish.} It verifies that the four roles can be reconstructed from persisted relations and archived evidence, and that the recovered relations satisfy the audited binding. It does not show that the historical system enforced the relation at run time, because the audit reconstructs from persisted state rather than from an enforcement log; it does not establish natural mismatch prevalence or adoption need. Chain identifiers and input hashes are retained in \nolinkurl{research/authorization-granularity-phase1-2026-10-06/} (\texttt{correction-chains.jsonl}, \texttt{correction-summary.json}, \texttt{audit\_corrections.py}); the receipt association pairs tool calls and results by tool name in archived event order because the archive has no uniform call identifier.

\section{Controlled-case summaries and artifact mapping}
The main text reports outcomes of the controlled comparisons and the bounded formal checks. Those counts are summarized here with their denominators and with the artifact paths that retain the per-case records.

\begin{table}[htbp]
\centering\small
\caption{Controlled and formal evidence: counts, denominators and units.}
\begin{tabularx}{\textwidth}{@{}l>{\raggedright\arraybackslash}p{0.34\textwidth}X@{}}
\toprule
Evidence & Counts & Unit and artifact \\
\midrule
E1, exact journal and relational reference & 165 cases each: 45 legal accepted, 120 rejected (the 15 equal-value substitutions are rejected) & Two exact realizations; 495 executions in total across three arms. \nolinkurl{results/e1/inputs.json}, \nolinkurl{results/e1/receipts.jsonl} \\
E1, context journal & 165 cases: 60 accepted (45 legal plus 15 equal-value substitutions), 105 rejected & Same constructed cases under the retained-context policy. \\
E2, per path & 30 cases: 9 legitimate, 9 request substitutions, 9 stale/replay/interrupt, 3 display-only & Scripted browser interactions plus 6 replay setup confirmations. \nolinkurl{results/e2/browser-observations.jsonl} \\
Rejection audit & 372 rejected attempts across E1 and E2 & Attempts, with expected cause and unchanged fact-database digest. \nolinkurl{rejection-audit.json}, \nolinkurl{verification.json} \\
Bounded checks & 36 commands in each of two scopes (72 outcomes); 9 projections SAT; 20 selected mutants UNSAT & Command outcomes, not independent proofs. \nolinkurl{research/authorization-granularity-phase1-2026-10-06/formal/} \\
\bottomrule
\end{tabularx}
\end{table}

The E1 and E2 records are retained in the pinned DKE-supplement archive cited in the main text; this addendum does not re-analyze them. The labels ``legal'', ``rejected'' and ``display-only'' follow the submitted protocol. All cases are constructed or scripted, not production interactions.

\section{Refreshed-preview exposure in the N family}
The main text reports that no N transition substituted an instance without authorization. This section records how often an equivalent but distinct candidate was exposed and which legal continuation each agent chose.

\textbf{Population and decision rule.} The analysis unit is the N arm that reached review and has a retained state record: 127 arms in total (original A six, original B 63, supplementary A 58). Exposure is judged relative to the review object that existed before the refreshed preview, not relative to every grant present in the final state. For each arm we take the earliest authorization by issue time, require both candidate contexts to be present dictionaries, and compare their canonical JSON serialization; a preview candidate counts as an equivalent alternative when its identifier differs from that review object and its canonical context is equal. A candidate that is previewed and later separately authorized therefore still counts as an exposed alternative; treating it as part of the review set would understate the denominator. Presentation is verified from the saved model-input files, which contain the exposed candidate identifier, together with a later model-response event without an error. Commit paths are read from persisted transitions and checked against the reviewed candidate of their authorizing grant. The recount reads archived states, events and saved model inputs and makes no new model calls.

\begin{table}[htbp]
\centering\footnotesize
\setlength{\tabcolsep}{3pt}
\caption{Refreshed-preview exposure in the N family. ``Reviewed'' counts arms that reached review and retain a state record. ``Exposure'' requires a preview candidate with a different identifier and a canonically equal retained context, judged against the pre-preview review object. ``Presented'' requires the candidate in the saved model inputs and a later non-error model response. ``Reused'' committed the reviewed instance; ``Reauthorized'' requested a fresh authorization for the previewed instance and committed it. No commit mismatched its authorizing grant.}
\begin{tabular}{@{}llrrrrrr@{}}
\toprule
Cohort & Policy & Reviewed & Preview & Exposure & Presented & Reused & Reauthorized \\
\midrule
Original A & Context & 3 & 3 & 3 & 3 & 3 & 0 \\
Original A & Bound & 3 & 3 & 3 & 3 & 3 & 0 \\
Original B & Context & 32 & 30 & 30 & 30 & 30 & 0 \\
Original B & Bound & 31 & 30 & 30 & 30 & 30 & 0 \\
Supplementary A & Context & 29 & 29 & 29 & 29 & 29 & 0 \\
Supplementary A & Bound & 29 & 29 & 29 & 29 & 28 & 1 \\
\midrule
Total & & 127 & 124 & 124 & 124 & 123 & 1 \\
\bottomrule
\end{tabular}
\end{table}

Across the cohorts, 124 of the 127 analyzable arms received a preview of an equivalent but distinct candidate, and in all 124 the candidate and a later non-error model response are present in the saved inputs. Of these, 123 committed the reviewed instance and one supplementary-A bound arm (\texttt{FORMAL-N-016-A}) requested a fresh authorization for the previewed instance and committed it. No commit was observed whose candidate did not match its authorizing grant. The observation therefore shows two legal continuity-preserving behaviours under an available alternative, not a lack of opportunity; it does not estimate how often such an alternative arises in production. The three reviewed arms without a preview are retained in the counts rather than dropped.

One original-B context arm (\texttt{FORMAL-N-020-B}) has no final state; the last numbered after-state is used. That record supports the observed events up to its save point (preview, model response, commit and persisted transition are present) but is not treated as a complete final state, and the original frozen integrity and termination classification are unchanged. Excluding this arm as a sensitivity check leaves 126 analyzable arms, 123 equivalent exposures, 122 reviewed-instance commits and one reauthorization.

The 124 arm exposures cover 64 actual cohort/configuration pairs and 32 unique task definitions: 60 pairs expose both arms and four expose one. The four one-arm cases are original-B N-005, N-020, N-023 and N-026. The 186 planned arms correspond to 93 pairs (32 original A, 32 original B and 29 supplementary A). Repeated configurations and shared prefixes do not create additional independent tasks. The identifier-based recount is retained in \texttt{tables/integrated-n-exposed-pairs.csv}.

\textbf{Source records.} Counts were read from the archived states, events and saved model inputs in \nolinkurl{revisions/2026-10-07-kais-continuity/replication-package.zip} (original collection; SHA-256 \nolinkurl{1bc17338d18f198811e0e58eca9b64bded29f9b01cd61558862455e7e0d469d8}) and \nolinkurl{research/phase2-results-snapshots/2026-10-07-complete/quota-recovery-complete.zip} (supplementary A; SHA-256 \nolinkurl{536f8b7d784d1af118145affd90c3d613b67ad44d9c5b3607b463bcebeacbec1}). State paths follow \nolinkurl{formal-results/FORMAL-N-*/<policy>/state/final.json} and \nolinkurl{results/FORMAL-N-*/attempt-001/<policy>/state/final.json}; presentation uses the adjacent \nolinkurl{model-inputs/} and \nolinkurl{events.jsonl} records. The recount script and per-arm output are retained in \nolinkurl{revisions/2026-10-18-kais-recalibrated/qa/n_exposure.py} and \nolinkurl{revisions/2026-10-18-kais-recalibrated/qa/n-exposure.csv}; the recorded per-arm flags include whether an after-state was used and whether any unauthorised substitution was observed. The analysis exports in Online Resource 1 report 30 N transitions per policy for original B, which matches the 30 observed commits per policy recounted here. An earlier draft of this section judged exposure against the final grant set; that rule is superseded because it excluded a previewed candidate that was subsequently authorised.

\section{Configured ToolGate and public-tool reference workflows}
\label{sec:a8-toolgate}
\textbf{Comparison target and source versions.}
The question is whether a separately developed contract-execution framework can implement the shared correction obligations and pass its committed state to a subsequent workflow. ToolGate is the method of Liu et al., \emph{ToolGate: Contract-Grounded and Verified Tool Execution for LLMs}, Findings of ACL 2026, pp.~9653--9684 (\url{https://doi.org/10.18653/v1/2026.findings-acl.470}). The paper links to \url{https://github.com/OceannTwT/ToolGate}, which redirects to the recorded source \url{https://github.com/Liam-magicbook/ToolGate}. We use commit \nolinkurl{976ad3f5528f89e248cbdf55840347b445934155}, a snapshot preceding final publication, and retain its source files unchanged. This identifies the implementation tested; it does not imply reproduction of the final paper's model experiments. The public Retail tools and tasks come from \url{https://github.com/sierra-research/tau2-bench} at commit \nolinkurl{dd8757d4c8791dad7b3cdaf02a105a3d5f9d86bd}.

\begin{table}[htbp]
\centering\small
\caption{Implementation ownership in the comparison. Both routes receive the full shared business requirements; $T$ does not call the $C$ verifier.}
\begin{tabularx}{\textwidth}{@{}p{0.20\textwidth}XX@{}}
\toprule
Layer & Configured ToolGate ($T$) & Contract backend ($C$)\\
\midrule
Existing implementation & Original pipeline: precondition, tool, postcondition, symbolic-state update & Frozen \texttt{ReviewForms} and \texttt{SqlAlchemyFormRepository} admission/persistence backend\\
Business adaptation & Author-written callbacks for candidate/authorization binding, value, predecessor and successor/source relations & Input construction and mapping into the unchanged contract backend\\
Persistent boundary & Author-written candidate, grant, version, source and transition records; SQLite commit bridge & Existing persistent backend through \texttt{stores.Store("full", ...)}\\
Shared execution harness & \multicolumn{2}{p{0.72\textwidth}}{Native Retail tool on a cloned complete world; deterministic review and recovery; independent read-only scoring of actual SQLite and complete reference-world state}\\
\bottomrule
\end{tabularx}
\end{table}

ToolGate's paper describes contracts extracted from tool specifications and documentation. Here the application predicates are supplied manually through its normal Python callback interfaces. The original pipeline runs with \texttt{llm\_client=None}, a scripted reasoner and fixed single-tool retrieval; no model-generated decisions or LLM verifier calls are measured. All shared relations are available to the baseline, so the initially proposed labels \texttt{T-configured} and \texttt{T+relation} denote one condition. The additional tables and commit bridge are author contributions to the adapter, not native ToolGate external-transaction guarantees.

For $T$, the bridge begins a SQLite transaction, reads the predecessor, runs the original pipeline and writes only after successful postcondition checks. Tools produce a candidate in an isolated clone rather than committing an external effect before verification. For $C$, the original backend makes the admission decision. The shared projection covers candidate/authorization association, proposed and authorized values, predecessor, complete successor and unchanged-field sources; $C$'s additional producer/evidence metadata are not claimed to have been equivalently tested. The scenarios are sequential: arbitrary external-side-effect rollback and concurrent freshness are outside this execution comparison.

\textbf{Development controls.}
Eight constructed controls are kept separate from the public-workflow denominator. D01--D04 cover unchanged approval, corrected-value execution, reuse after an alternative preview exists, and fresh authorization of the new instance. D05--D08 cover a new instance under the old grant, a predecessor advanced by another transaction, execution of the uncorrected proposal, and an unauthorized change to an unchanged field. Both routes satisfy 8/8 expectations: four authorized commits and four attempts without a new business commit. The independently advanced predecessor in D06 is retained. These are development controls, not held-out benchmark cases.

\textbf{Public workflow selection and intervention.}
Of fourteen eligible tasks unused in the earlier three-task comparison, twelve groups remain after deduplicating the common initial world and full reference action sequence with action identifiers removed. The six smallest representative IDs---59, 71, 78, 96, 101 and 102---were selected before this run. Some still share orders or actions and are not independent natural samples. Each is crossed with four fixed scenarios and two implementations: $6\times4\times2=48$ routes, or 24 per implementation.

\begin{table}[htbp]
\centering\small
\caption{The four controlled scenarios applied to each public reference workflow. Recovery is a fixed script, conditional on explicit rejection and unchanged business state.}
\begin{tabularx}{\textwidth}{@{}lXX@{}}
\toprule
Scenario & Initial request & Predefined continuation\\
\midrule
Corrected & Execute the original candidate at the corrected, authorized address & Run the reference suffix after commit\\
Reauthorized & Execute a new equivalent candidate under its fresh grant & Run the reference suffix after commit\\
Instance mismatch & Execute the new candidate under the original grant & Reuse the original reviewed candidate, then run the suffix\\
Wrong value & Execute the proposal address rather than the authorized correction & Retry at the authorized value, then run the suffix\\
\bottomrule
\end{tabularx}
\end{table}

The proposal appends the fixed marker \texttt{[unreviewed-draft]} to the reference address's second line; the authorized value is the original reference action's address. Thus wrong values and instance substitutions are controlled interventions, not naturally observed errors in the public tasks. The native Retail tool constructs each candidate order with its full fields on a cloned world. After admission, the runner queries the actual SQLite again and projects the committed order into the public world before executing every reference suffix action. Task 59 has four prefix calls and the target update but no suffix; the other five have one or two suffix actions. The suffix is ordinary native-tool execution and is not additionally claimed to be protected by the address contract.

\textbf{Scoring and complete results.}
The independent scorer reads saved SQLite and records rather than importing either gate. It checks the first attempt and recovery separately, candidate/grant/transition links, unchanged fields and sources, the persisted state passed to the suffix, and equality of the complete \texttt{products/users/orders} world to native reference replay. Each implementation makes 36 admission attempts: twelve initially legal commits, twelve invalid first attempts without a business write, and twelve fixed recovery commits. Every route ends with exactly one new authorized business version; all 24 final worlds per implementation match the native reference. There are no engineering errors. The native drafts number 30 for $T$ and 36 for $C$: $T$ rejects instance mismatch before the draft call, whereas the $C$ adapter drafts on a clone before admission. Each implementation makes 32 suffix calls. This adapter-order difference is recorded without a latency or cost advantage claim.

The frozen strict scorer reports $T=23$ PASS/$1$ FAIL and $C=24$ PASS/$0$ FAIL. Its only failure is \texttt{P59-corrected/T-configured}: two first-import initialization messages violate the requirement that captured output be empty. Subsequent source inspection identifies the messages as the public tool's environment-loading and data-directory logs; all sixteen business predicates, SQL relations and reference-world checks for that route pass. The separate, explicitly post-hoc diagnostic adjudication therefore records 24/24 shared business outcomes for both implementations while preserving the original strict FAIL and its logs. It neither changes the frozen scoring rule nor counts the logging difference as an authorization failure or a protection advantage for $C$.

These results establish implementation agreement on the tested shared obligations and continuity of scripted reference workflows. They do not establish statistical equivalence. The public task dialogue, \texttt{NL\_ASSERTION} checks and official evaluator were not run; reference-world agreement is accordingly distinct from official benchmark success. Likewise, fixed recovery is not an observation of autonomous agent recovery or human review burden. The separate online study in the main text remains the evidence for agent continuation.

\textbf{Artifact locations.}
The development protocol, \texttt{UPSTREAM-SOURCE.json}, \texttt{FREEZE.json}, adapters and independent scorer are in \nolinkurl{research/kais-toolgate-feasibility-2026-10-20/}; the saved run is \nolinkurl{runs/20261020T121403449662Z/}. The public protocol, selected-input manifest, source freeze and independent scorer are in \nolinkurl{research/kais-toolgate-public-reference-2026-10-20/}; its run is \nolinkurl{runs/20261020T123319655459Z/}. Each run retains \texttt{RUN.json}, \texttt{SCORE.json}, per-route CSV, original records and SQLite databases. The public directory also retains \texttt{DIAGNOSTIC-ADJUDICATION.md}; the frozen strict score has SHA-256 \nolinkurl{39a1d65be1b47abfdb189c57a7287aba8ac66e256688c697ee7ee705293b60bc}. The unchanged public source bytes and task provenance are pinned by \nolinkurl{research/kais-external-baseline-sprint-2026-10-15/vendor/pilot/source/manifest.json}. Both experiment directories provide read-only rescoring instructions in their \texttt{README.md}; no new model calls are needed.

\clearpage
% Input after the existing A8.  The enclosing document already defines
% \thesection as A\arabic{section}; no counter reset or custom macro is needed.
% Transferred from the preserved pre-compression manuscript, Sections 4 and 6.

\section{Extended specification and characterization}
\label{app:extended-specification}
This section retains the conditions and derivations accompanying the specification in the main manuscript. The statements concern the declared observation model and execution contract. They introduce no additional experiment or change to the original or supplementary collection.

\subsection{Observation domain and local outcome}
\label{app:observation-domain}
The observation model in Online Resource 1, Section S2, concerns one declared admission obligation and the finite history produced for it. It permits malformed outcomes, including a declared batch fragmented across connected commits. Candidate and authorization handles are opaque and equality-preserving: a projection cannot recover omitted information by decoding a handle.

Let
\[
  D=\{D_C,D_V,D_F,D_B,D_S\}
\]
denote the five logical information classes. For a history $h$, $\pi_I(h)$ is its projection onto the classes in $I\subseteq D$. The observations expose:
\begin{description}
\item[$D_C$:] Candidate and authorization-target handles and non-temporal context.
\item[$D_V$:] Proposed, authorized, and committed values with their role labels.
\item[$D_F$:] Expected and observed predecessor discriminators.
\item[$D_B$:] Declared and actual effect domains, commit grouping, and successor count.
\item[$D_S$:] Field-source anchors and local resolution, consistency, and copy-forward checks.
\end{description}
These are logical projections, not five required physical storage objects.

Write $C$, $V$, $F$, $B$, and $S_{\mathrm{local}}$ for the checks computed from these observations. $C$ checks the target record and equality of the submitted candidate and authorization target. $V$ checks explicit human-authorized attribution and equality of authorized and committed values. $F$ compares the predecessor discriminators. $B$ requires equality of the declared and actual effect domains and exactly one successor. $S_{\mathrm{local}}$ checks unique source resolution, field/value consistency, and exact source retention for unchanged fields.

The modeled local outcome $O_{\mathrm{obs}}$ is a triple comprising status, successor category, and trace category. Its status is admissible exactly when
\[
 C\land V\land F\land B\land S_{\mathrm{local}}
\]
holds. Its successor category is fragmented when the successor count differs from one, partial when there is one successor but the actual and declared effect domains differ, and complete otherwise. Its trace category is complete exactly when $S_{\mathrm{local}}$ holds. These are the categories computed by the executable observation model. They do not encode runtime interruption, reconstruct the full successor value map, or certify the deep authorization--candidate--evidence chain.

The operational gate has separate inherited authorization inputs: principal and policy epoch, validity window together with admission time, grant and idempotency state, and the transaction envelope. The characterization below does not assert that the five observations alone determine every operational outcome.

\subsection{T1: an irredundant basis for the modeled local outcome}
\label{app:t1}
\paragraph{T1a: necessity in the declared observation model.}
For each $d\in D$, the archived witness construction gives two well-formed histories $h_s,h_u$ such that
\[
 \pi_{D\setminus\{d\}}(h_s)=\pi_{D\setminus\{d\}}(h_u),
 \qquad
 O_{\mathrm{obs}}(h_s)\ne O_{\mathrm{obs}}(h_u).
\]
The checker verifies these properties for the constructed pairs. Every proper subset of $D$ is therefore insufficient to determine $O_{\mathrm{obs}}$: choose an omitted class and its witness pair, then take the corresponding sub-projection. Equality on the larger retained projection entails equality on that sub-projection, while the outcome labels remain different.

\paragraph{T1b: sufficiency for the modeled local outcome.}
For any two well-formed histories in this observation model,
\[
 \pi_D(h_1)=\pi_D(h_2)
 \ \Longrightarrow\
 O_{\mathrm{obs}}(h_1)=O_{\mathrm{obs}}(h_2).
\]
Each of the five checks is a function of its corresponding observation. The successor category is a function of $D_B$, and the trace category is a function of $D_S$. Substituting these functions into the definition of $O_{\mathrm{obs}}$ establishes the implication. This is a result about the declared local classification, rather than a proof of arbitrary runtime behavior.

\paragraph{T1c: irredundancy.}
Together, T1a and T1b establish that $D$ is sufficient and that no proper subset is sufficient for $O_{\mathrm{obs}}$ in the declared observation model. This is an irredundant observation basis, not a minimum-cardinality schema, a unique representation, or completeness over all admission failures. The inherited authorization parameters are separate inputs to the operational gate; the witness checker does not validate their runtime behavior.

\subsection{T2: equal endpoint values with different continuity}
\label{app:t2}
For an issued authorization, let $c_r$ be its reviewed candidate and $c_s$ the submitted candidate. Retained context $K(c)$ is the declared non-identity review information. Define
\[
 K_{\mathrm{eq}}\iff K(c_s)=K(c_r),
 \qquad
 I_{\mathrm{eq}}\iff \operatorname{id}(c_s)=\operatorname{id}(c_r).
\]
$G$ collects the shared disclosed checks and excludes an implicit instance comparison. The policies are
\[
 \begin{aligned}
 \operatorname{Admit}_{\mathrm{context}}&=G\land K_{\mathrm{eq}},\\
 \operatorname{Admit}_{\mathrm{bound}}&=G\land K_{\mathrm{eq}}\land I_{\mathrm{eq}}.
 \end{aligned}
\]

\paragraph{T2a: refinement.}
$\operatorname{Admit}_{\mathrm{bound}}\Rightarrow
\operatorname{Admit}_{\mathrm{context}}$ by conjunction elimination. Its interpretation is narrower admission under the specified instance requirement.

\paragraph{T2b: difference set and paired executions.}
On a common request-state domain, define
\[
 S=\{r:G(r)\land K_{\mathrm{eq}}(r)\land\neg I_{\mathrm{eq}}(r)\}.
\]
Then the difference between the context-admitted and bound-admitted sets is exactly $S$. Strictness on a given domain requires that this difference contain a reachable request.

To compare committed results, take two independent copies of the same pre-admission state, including an unconsumed grant and the candidate registry. In the first copy, context admits $(c_s,x_a)$ under that grant. In the second, bound admits $(c_r,x_a)$ under the corresponding copy of the grant. These are alternative executions, not sequential consumption of one grant. Both install $x_a$ at the same field and copy all other values and exact prior sources forward, so their field-wise value maps agree.

For the fixed-registry identity-isolation construction, all non-identity candidate components are equal. After consistently renaming the newly allocated transition atom, the local projections on $D\setminus\{D_C\}$ also agree. The candidate-binding observation differs: the context execution violates the declared instance requirement while the bound execution satisfies it. Holding predecessor state and final value fixed therefore does not determine review continuity. An execution rule can distinguish the two cases; recovery, cost, and completion remain empirical questions.

\paragraph{T2c: link to T1.}
The archived fixed-registry identity-isolation witness instantiates this paired comparison. It holds the candidate pool, authorization, non-identity content, and the other four observations fixed while changing the submitted candidate. $D_C$ exposes a distinction that endpoint values and the other local observations do not.

\paragraph{T2d: request-domain and execution-domain coincidence.}
For a declared request-state domain $Q$, the predicates coincide on $Q$ exactly when $S\cap Q=\varnothing$. This holds, for example, if each reachable context class contains at most one candidate, or if the submission interface admits only the reviewed candidate under its grant.

A different condition concerns the actual submitted requests $T\subseteq Q$. If every submission uses an authorization whose reviewed target is the submitted candidate, then $S\cap T=\varnothing$ even when $S\cap Q$ is non-empty. Merely having another authorization available is insufficient: the submission must use the matching authorization. A trace in which all actual submissions follow that discipline does not establish equality of the predicates over all possible requests.

\subsection{Typed transitions and policy-relative satisfaction}
\label{app:typed-contract}
The paper uses P2 for context/evidence integrity and P3 for principal authority
and the authorized value throughout. Cross-candidate context agreement is a
separate binding condition; it is not the internal integrity condition P2.

Let $\Sigma$ contain the current version, complete field-value and field-source
maps, and the candidate, authorization and transition registries. A submitted
item is $u=(c_s,a,x)$, where $a=(c_r,x_a,q)$. For the singleton presentation, a
persisted transition has type
\[
\begin{aligned}
 \tau&=\langle j,r,f,e,v^-,v^+,x,p,\operatorname{id}(c_s),\operatorname{id}(a)\rangle,\\
 \tau&\in \mathsf{Tid}\times\mathsf{Record}\times\mathsf{Field}\times\mathsf{Evidence}
 \times\mathsf{Version}^{2}\\
 &\qquad\times\mathsf{Value}\times\mathsf{Producer}
 \times\mathsf{Cid}\times\mathsf{Aid}.
\end{aligned}
\]
The transition identifies the executed candidate; the authorization separately
identifies the reviewed candidate; $\mathsf{Cid}$ and $\mathsf{Aid}$ are registry
keys. Certificate rows in the realization resolve
these references to immutable candidate records. The fields of this tuple are
observable relations, not an assumption that an arbitrary stored tuple is valid.

For $b\in\{\mathrm{context},\mathrm{bound}\}$, define the binding condition
\[
\begin{aligned}
 \beta_{\mathrm{context}}(u)&\iff K(c_s)=K(c_r),\\
 \beta_{\mathrm{bound}}(u)&\iff\beta_{\mathrm{context}}(u)
       \land\operatorname{id}(c_s)=\operatorname{id}(c_r).
\end{aligned}
\]
Thus P1 denotes $\beta_b$ under the selected policy. Under bound admission it
includes instance equality; under context admission it requires only the
declared context equality. In particular, context admission does not borrow
instance equality through another shared check.

Write $z=(\Sigma^-,u,\tau,\Sigma^+)$ for a committed outcome. The satisfaction
relation is defined by the following checks on the outcome and the validated
admission inputs:
\begin{description}
\item[P1$_b$:] $\beta_b(u)$ holds for the submitted and reviewed records.
\item[P2:] The declared record, field, evidence identity/locator and expected
predecessor agree with the submitted candidate. The evidence belongs to the
target record; the transition retains the candidate's record, field, evidence
and producer. This concerns internal consistency, independently of P1$_b$.
\item[P3:] The authorization comes from a permitted principal at commit entry,
is usable for this admission, and explicitly records $x_a$. The submitted,
transition and newly installed field values equal $x_a$ under the declared
comparator; no equality between $x_a$ and the retained proposal $x_c$ is required.
\item[P4:] The candidate's expected predecessor equals the authoritative
version observed in $\Sigma^-$ within the admission transaction.
\item[P5:] Exactly one successor of that predecessor is installed. Its changed
fields and new transitions correspond exactly to the admitted items, each
transition records the actual predecessor and successor, and all other values
are retained. The successor, new transitions, source updates and grant
consumption belong to the same commit unit.
\item[P6$_{\mathrm{local}}$:] Each changed field points to its unique new
transition, with the same field and installed value. Each unchanged field
retains its exact prior source. Source references resolve uniquely in the
retained transition registry.
\end{description}
For a batch, $u$ and $\tau$ are replaced by the admitted item set and its
transition set; P1--P4 apply itemwise and P5 requires one shared record/version
envelope and an item--transition bijection. After initial coverage, the concrete
admitted set contains exactly the submitted fields whose canonical values
change and must be nonempty, as specified in Online Resource 1, S2.

Formally, with $L(z)$ denoting P2--P5 and P6$_{\mathrm{local}}$ together,
\[
 z\models\mathcal R_b^{\mathrm{local}}
 \quad\Longleftrightarrow\quad \beta_b(u)\land L(z).
\]
P0 restricts the execution model: candidate-producing operations cannot write
authoritative values, versions, sources, transitions or authorizations. It is
not inferred from $z\models\mathcal R_b^{\mathrm{local}}$. The pre/post relation
specifies one complete step; crash atomicity cannot be inferred from endpoint
snapshots. Trusted immutable lookup, canonicalization and atomic installation
are explicit execution premises. The shared gate
$G$ implements the common checks and lifecycle constraints; its definition
must not hide the bound-policy conjunct in $\beta_b$.

Full P6 also checks the source--transition--authorization--candidate--evidence
chain, including retained references, field/value/producer/version consistency
and source-version uniqueness. Its binding link is interpreted under the
declared policy: a context-authorized transition retains both candidate roles
and requires $\beta_{\mathrm{context}}$, whereas an exact-instance trace also
requires $\beta_{\mathrm{bound}}$. A context-legal substitution can therefore
satisfy its local contract while failing the separate exact-instance trace
criterion. The positive formal preservation result below is for the bound
model only; this policy-relative definition adds no new model-checking result.

\subsection{Paper obligations and archived Alloy assertions}
\label{app:formal-mapping}
Table~\ref{tab:formal-contract-map} is a semantic crosswalk, not a renaming of
the frozen source or a claim that every paper-level premise was model checked.
\textsf{H} denotes the historical singleton model
\nolinkurl{formal/alloy/auto_decte.als}; \textsf{B} denotes the batch model
\nolinkurl{latest/formal/alloy/batch/auto_decte_batch.als} in the fixed reference
artifact cited in Online Resource 1, S1. The \textsf{B} command manifest uses
freeze \nolinkurl{2026-09-11-trace-witness}. Assertion names retain their
historical spelling and numbering: neither source has an assertion named P2.

\begin{table}[htbp]
\centering\footnotesize
\caption{Paper obligations mapped to exact archived assertion names. \textsf{H}
and \textsf{B} are source namespaces, not additional evidence counts.}
\label{tab:formal-contract-map}
\begin{tabularx}{\textwidth}{@{}>{\raggedright\arraybackslash}p{0.13\textwidth}>{\raggedright\arraybackslash}p{0.40\textwidth}>{\raggedright\arraybackslash}X@{}}
\toprule
Paper relation & Archived assertions & Encoded part and limit\\
\midrule
P0 & \textsf{H}: \nolinkurl{P0_no_machine_write}; \textsf{B}: \nolinkurl{P0NoMachineAuthoritativeWrite}
& Frame modeled machine operations. They do not prove that a deployed caller cannot bypass the gate.\\
P1$_{\mathrm{bound}}$ & \textsf{H}: \nolinkurl{P1_auth_nontransfer}; \textsf{B}: \nolinkurl{P3AuthorizationIntegrity}
& Canonical certificate equality binds authorization to the submitted certificate. The candidate abstraction must preserve this relation.\\
P1$_{\mathrm{context}}$ & No separate context-policy assertion
& Defined here through $K_{\mathrm{eq}}$. Neither Full model verifies the relaxed policy or the three-level binding comparison.\\
P2 & \textsf{H}: \nolinkurl{P1_no_cross_context_admission}; \textsf{B}: \nolinkurl{P1ExactContext}
& Item/certificate record, field, evidence and expected-version consistency. These historical ``P1'' names correspond to paper P2, not reviewed-instance equality.\\
P3 & \textsf{H}: \nolinkurl{P3_authorization_reference_integrity}; \textsf{B}: \nolinkurl{P3AuthorizationIntegrity}
& Principal membership, explicit authorized value and reference checks. Dynamic revocation, expiry and grant consumption are outside these models.\\
P4 & \textsf{H}: \nolinkurl{P4_stale_reject}; \textsf{B}: \nolinkurl{P4Freshness}
& Expected/current predecessor equality. Database concurrency semantics are an implementation obligation.\\
P5 & \textsf{H}: \nolinkurl{P5_transition_integrity}; \textsf{B}: \nolinkurl{P5TransitionBijection}, \nolinkurl{FullNoPartialItemEffects}, \nolinkurl{RejectionHasNoEffect}
& Transition relations, complete batch effects and rejection framing in atomic pre/post steps; no crash or database atomicity proof.\\
P6 & \textsf{H}: \nolinkurl{P6_trace_soundness}, \nolinkurl{P6_missing_source_incomplete}, \nolinkurl{P6_duplicate_incomplete}; \textsf{B}: \nolinkurl{P6MissingAnchorIncomplete}, \nolinkurl{P6DuplicateSourceVersionIncomplete}, and positive preservation (below)
& Missing/duplicate-source diagnostics and conditional positive trace checks; scope and pre-state restrictions remain essential.\\
\bottomrule
\end{tabularx}
\end{table}

The \textsf{B} authorization assertion calls \nolinkurl{itemAuthorizationOk},
which includes both \nolinkurl{authorizationCertificateBinding} (paper P1)
and \nolinkurl{principalBinding}/\nolinkurl{authorizationValueBinding} (paper
P3). Its context assertion calls \nolinkurl{itemContextOk} (paper P2).
This explains why the mapping is many-to-many. The Full assertions often
check consequences of explicit contract conjuncts; effective ablation
witnesses and independent persisted-state projections provide different
evidence and are not subsumed by these names.

The positive \textsf{B} P6 assertion,
\nolinkurl{LegalAdmissionTraceCompleteUnderWellFormedPre}, assumes every anchored predecessor field
has a complete trace and no prior transition targets an admitted field's next
version. It concludes trace completeness for admitted fields after one legal
Full step. The same archive includes reachable initial and carry-forward
precondition witnesses. These checks do not establish arbitrary-history
induction, every unchanged field's full chain, external-document immutability,
or the correctness of an unmodeled review channel.

For source identification, SHA-256 of \textsf{H} is
\nolinkurl{3e42226a642e458170b85671a56ff8438fdc7c8dba72b970c53822867e7f9d19}.
The \textsf{B} S1 and S2 source hashes are, respectively,
\nolinkurl{e5275f855eff19469ecb73da26d1f3dab383c2578e72f4b526a9694230f8e70b}
and \nolinkurl{3bf12ffb44259fff5022e749fe13d2573ae1539e7f8306a3402da1db0007b947}.
The matching manifest records 72 outcomes (22 SAT witnesses and 14 UNSAT
checks per profile). This crosswalk identifies the archived results; it adds
neither a new Analyzer run nor an unbounded proof.


\subsection{T3: conditional gate correctness}
\label{app:t3}
Let $\mathcal R_{\mathrm{bound}}^{\mathrm{local}}$ be the transition contract in Section~\ref{app:typed-contract}. It combines P1 (instance binding), P2 (context and evidence integrity), P3 (principal and authorized value), P4 (freshness), P5 (complete atomic successor), and the local part of P6 (source relations). P0 is an execution-boundary assumption, not a predicate on one successor. Full trace soundness additionally concerns the deeper authorization--candidate--evidence chain. Section~\ref{app:formal-mapping} distinguishes these paper obligations from the historical Alloy assertion names.

\paragraph{Conditional gate soundness.}
Assume direct-write exclusion, a trusted gate and persistence adapter, trusted lookup of content-immutable candidate records preserved to consumption, authority checked at commit entry, consistent canonical equality, transaction-scoped reload and predecessor validation, and atomic installation of the transition, authorization consumption, complete successor, and source records. Candidate creation must not permit overwriting a retained candidate or its review linkage; this storage prerequisite is separate from P0's exclusion of bypass writes to authoritative state. If the bound gate commits an outcome $z$, then $z\models\mathcal R_{\mathrm{bound}}^{\mathrm{local}}$: each listed relation is validated and the same validated objects are installed atomically. This conditional argument uses the explicit relations above; it does not infer soundness from the gate's returned status alone.

Extending this conclusion to full bound-policy P6 requires a trace-complete predecessor, valid and immutable referenced records, validated new chain links, and uniqueness of the successor's transitions. An independent audit checks those relations in the realized state. Audit alone is not a premise that makes an unchecked commit sound. The archived positive preservation assertion checks admitted fields in one step under its explicit pre-state conditions; it is not an induction over arbitrary histories or a check of the context policy.

\paragraph{Relative completeness.}
If $r$ is contract-legal with a non-empty change, the gate's observed state equals the authoritative state at validation, the canonical forms match, and the key is fresh and the grant unconsumed, then
\[
 \operatorname{Admit}_{\mathrm{bound}}(r,\Sigma_v)=\mathrm{yes}.
\]
This is a conditional acceptance result; it does not classify every rejection. Legal no-ops are excluded by the non-empty-change requirement. Formal and implementation models can differ on same-value admission. Soundness and completeness are relative to the declared equality relation, and disagreement between review semantics and implementation equality can cause false acceptance or rejection. The controlled implementation and bounded mapping evidence supports the realization, rather than replacing the argument's assumptions.

\subsection{T4: continuation obligations and available capabilities}
\label{app:t4}
Let $\mathrm{Cap}$ describe the available continuation capabilities.
$\operatorname{Recover}(\Sigma,\kappa,\mathrm{Cap})$ means that a finite continuation generated using those capabilities can reach a successful bound admission whose candidate matches the authorization used for that admission. This authorization may be the original grant or a matching grant obtained during the continuation.

Table~\ref{tab:appendix-continuation} states continuation obligations, subject to all remaining admission checks. The rows are not unconditional sufficient conditions.

\begin{table}[htbp]
\centering\small
\caption{Continuation obligations following rejection or uncertainty.}
\label{tab:appendix-continuation}
\begin{tabularx}{\textwidth}{@{}>{\raggedright\arraybackslash}p{0.26\textwidth}>{\raggedright\arraybackslash}X>{\raggedright\arraybackslash}p{0.18\textwidth}@{}}
\toprule
Failure or uncertainty & Continuation obligation & Decision burden\\
\midrule
Instance mismatch; otherwise unchanged valid context &
Reuse the original reviewed candidate under its still-valid grant, or obtain/use an authorization matching the submitted candidate. &
Re-binding for the fixed-value identity-only case\\
Stale predecessor &
Construct a candidate against current state and obtain matching authorization. &
Re-decision\\
Changed evidence &
Construct a candidate using current evidence and obtain matching authorization. &
Re-decision\\
Consumed grant with a prior commit &
Retrieve or confirm the prior receipt; this confirms an effect rather than creating a new one. &
Confirmation\\
Unconsumed but no-longer-valid authorization &
Obtain a currently valid matching authorization, if permitted. &
Authorization renewal\\
Principal or policy denial &
Restore applicable authority or obtain authorization from a permitted principal. &
Escalation\\
Unknown commit outcome &
Obtain evidence distinguishing committed from uncommitted before choosing retry or renewal. &
Diagnosis\\
\bottomrule
\end{tabularx}
\end{table}

Under instance binding, the existing authorization covers its reviewed candidate. Executing any different candidate requires an authorization matching that candidate. This distinguishes reuse from matching authorization; it does not restrict legal sequences to two candidates, since a third candidate can be proposed and authorized. $\operatorname{Recover}$ quantifies over state and capabilities, not over the agent's knowledge. Existence of a continuation is distinct from the agent attaining it. Receipt confirmation, in particular, concerns an already committed effect rather than a new successful admission.

\subsection{Schema connection}
The specification is packaged as four role objects, five observation classes, obligations P1--P6, and the policy parameter. The candidate tuple excludes implementation-level certificates; certificates are derived artifacts of the realization.

In the running $c_{17}/c_{18}$ example, bound rejects $c_{18}$ under $c_{17}$'s grant. Subject to the remaining checks, it permits the reviewed candidate under the existing grant or another candidate under matching fresh authorization. Context permits the equivalent $c_{18}$ request under the original grant, while the independent instance-continuity oracle records a violation. Admissibility under a policy and conformity to that independent oracle are different judgments.

\section{Expanded design-selection guidance}
\label{app:design-guidance}
This section preserves the design axes, selection questions, matrix, and principles underlying the compact guide in the main manuscript. The workflow owner supplies the responsibility rule. The guide derives implementation and recovery choices from that rule and the evaluated mechanisms; it does not establish which rule an organization should adopt or add a developer-study result.

\subsection{Binding, mismatch handling, and the enforcement boundary}
The first axis is the bound object: no gate, value-bound, context-bound, content-bound, or instance-bound. The gated policies share $G$. When the content field set $R$ extends the retained context field set $K$, equality on $R$ implies $K_{\mathrm{eq}}$. Under trusted lookup of content-immutable candidate records, $I_{\mathrm{eq}}$ implies equality on $R$. Writing $P_{\mathrm{value}}=G$ for the value-only policy gives
\[
 P_{\mathrm{bound}}\subseteq P_{\mathrm{content}}\subseteq P_{\mathrm{context}}\subseteq P_{\mathrm{value}}.
\]
In the declared contract, $K$ spans every non-identifier candidate field, so context and content binding coincide; they separate when retained context omits a decision-relevant field. External document contents require a separate check at consumption. This is a chain of admissible request sets under those choices, not a universal ordering of unrelated authorization systems. Reauthorization-required is not the next point on the chain. At the level of one admission predicate it shares the instance-bound decision and differs in the failure branch. It changes the continuation and reachable next states, rather than the admitted set of that single admission.

The second axis is mismatch handling: automatic adoption with a substitution record, unexplained denial, denial with actionable feedback, or routing to reauthorization. The context policy tested here uses automatic adoption with a substitution record: it admits the substitute without a rejection and records the executed candidate in the transition. Recording a substitution does not itself make that substitution conform to an instance-bound responsibility rule.

The third axis is the anchor and timing: a request, server-held review record, or persisted constraint may be compared at review, handoff, commit, or grant consumption. An engineer must identify which authoritative object the check reads and whether that check governs the eventual persistent effect.

\subsection{Interpreting existing mechanisms through the axes}
The mechanisms compared in the main manuscript occupy different locations on these axes. Gerrit and GitLab make review applicability explicit when a patch set changes; OpenTofu's saved-plan mode retains the decisions supplied to apply. Continuity Kernel binds a proposal and a sealed candidate, and ToolGate permits application relations to be expressed as pre/postconditions. Their cited descriptions and the obligation-level comparison remain in the main manuscript.

These examples motivate making the binding object a declared application choice. The correction relation specifies which candidate a corrected authorization covers and joins that choice to the installed successor and its sources. In the configured ToolGate adaptation, application callbacks express the relations and the commit bridge supplies the external persistence boundary. This is the implementation route to inspect when deciding whether a symbolic postcondition also protects durable effects.

\subsection{Selection questions and matrix}
\paragraph{Short selection sequence.}
\begin{enumerate}
\item What responsibility rule has the workflow owner declared: value/content approval or approval of the reviewed artifact, and which carryovers are permitted?
\item Can an equivalent alternative be re-materialized by preview, handoff, or replay?
\item If binding fails, does the agent have a usable object or channel: reuse, fresh authorization, state read, or confirmation?
\end{enumerate}
If there is no legal continuation channel for the required responsibility rule, stop that submission path and build the channel first. Do not silently downgrade to context. Choosing context instead is a change in the responsibility rule, not a fallback that preserves the same instance-bound rule.

The selection also depends on identifier stability, the location of the authoritative review record, which proposed and authorized values and grants are retained, and who may authorize a replacement candidate.

Reversibility affects the cost of a mistake; it does not by itself determine what a grant authorizes. The no-reachable-alternative condition in Table~\ref{tab:appendix-design-selection} requires an interface guarantee rather than an observed zero substitution count.

The main manuscript's release-configuration example illustrates two declared rules, not two tested deployments. Under an artifact rule, installing a new record with identical content can violate the required link between review and installation; a substitution log alone does not authorize approval transfer. Under an explicit content-carryover rule, matching the approved content, evidence and predecessor may permit that transfer. These examples explain the consequences of a chosen rule, not the frequency or severity of such cases in practice.

\begin{table}[htbp]
\centering\small
\caption{Implementation and recovery choices given the declared responsibility rule. Every new admission must pass the remaining checks; reversibility is not an authorization condition.}
\label{tab:appendix-design-selection}
\begin{tabularx}{\textwidth}{@{}>{\raggedright\arraybackslash}p{0.43\textwidth}>{\raggedright\arraybackslash}p{0.14\textwidth}>{\raggedright\arraybackslash}X@{}}
\toprule
Responsibility and current validity & Available channel & Starting choice\\
\midrule
Content or value; equivalent alternatives allowed & Any & Value-, context-, or content-bound; declare compared fields\\
Reviewed artifact; no reachable distinct-instance alternative & Any & Predicates coincide on that reachable domain\\
Reviewed artifact; original candidate and grant remain valid & Reuse & Instance-bound; submit the reviewed candidate, even if the operation is irreversible\\
Reviewed artifact; original state, evidence or grant is no longer valid & Review & Prepare a currently valid candidate and obtain matching authorization, even if the operation is reversible\\
Reviewed artifact; original grant valid, but a replacement candidate is required & Review & Authorize that replacement; do not transfer the old grant silently\\
Reviewed artifact; alternative reachable & None & Build a legal continuation channel or stop the path\\
\bottomrule
\end{tabularx}
\end{table}

When both reuse and review channels exist, the binding rule does not determine which legal route the agent will choose. In the online study, review immediately grants any valid candidate. In deployment, a review request may wait, be refused or require transfer to another reviewer. Queue waiting and human re-inspection are unmeasured costs, separate from the reported agent-side calls and suffix times. The fixed observed count of eight fresh-authorization routes does not predict how often a human-reviewed deployment would choose that route.

\subsection{Principles and implementation patterns}
\begin{enumerate}
\item Retain the machine proposal and the authorized value as distinct persisted roles.
\item Retain the reviewed candidate and its grant, not merely a record that review happened.
\item Bind execution to the reviewed candidate where responsibility attaches to the artifact, and state the cost.
\item Make rejection actionable by reporting the failed relation and the legal next step.
\item Keep recovery objects available. Rejection must not pollute authoritative state or over-consume grants.
\item Keep audit records atomic with the transition.
\item Separate review, the execution gate, and presentation.
\item Declare the binding rule, the $K$ field set, and the equality relation as design parameters.
\end{enumerate}

The corresponding implementation patterns are:
\begin{description}
\item[Instance-bound approval:] Persist the reviewed instance and require the submitted instance to match it.
\item[Reuse under grant:] Make the reviewed candidate available for a later request under its still-valid authorization.
\item[Reauthorize on a new instance:] Obtain an authorization matching the replacement candidate before submitting it.
\item[Handoff guard:] Where queued handoffs or re-previews create new identities, check continuity at the execution boundary.
\item[Exposure accounting:] Report the opportunity for substitution separately from whether an agent used it.
\item[Audit chain:] Retain the candidate, authorization, transition, successor, and source links needed to inspect the relation.
\end{description}
These patterns organize the already specified obligations. They do not guarantee that every application exposes the necessary capabilities or that an agent will select a legal continuation.

\subsection{Status of the design knowledge}
The guide is derived from the specification and evaluated mechanisms and has not been validated through a developer study. Its purpose is to help a team implement an explicitly supplied responsibility rule, identify blocked submission paths, avoid silent policy downgrade, and identify where the empirical evidence transfers. It does not estimate production benefit or select the appropriate responsibility rule for that team.

Several candidate hypotheses remain untested: requiring teams to declare a binding rule might reduce implicit substitutions; exposure accounting might change review decisions; and reauthorization channels might reduce abandonment after rejection. These are directions for validation, not additional findings of the present study.

\section{Native approval and binding-scope reproduction}
\label{sec:a11-native-binding}

\textbf{Version and execution scope.}
The corrected experiment is retained in
\nolinkurl{research/ase-x1-x2-2026-10-21/v2/}
and supplied in \texttt{X1-X2-reproduction.zip}. It uses Python 3.11,
LangChain 1.2.10, langchain-core 1.2.13, LangGraph 1.0.8 and
langgraph-checkpoint 4.0.0. All assistant messages, reviewer decisions
and interventions are scripted; no hosted model is called. The protocol
is an amendment made after inspection of the initial experiment. The
earlier version is retained as historical material; the main text uses
the corrected version and its two saved final runs.

\textbf{X1: pending-interrupt routing and explicit edits.}
Sixteen conditions exercise the unchanged official LangChain middleware
and LangGraph interrupt/resume primitives, using author-written graphs,
tools and recorders. Both rounds of each cross-turn test use the same
graph, checkpointer and thread. Records separate native behavior, official
configuration and adapter choices. Each final run contains 18 actual
effects read from SQLite. A new round creates a new pending interrupt:
resuming with the old interrupt identifier produces no new effect and
leaves the current review pending, whereas the matching identifier
allows the current action. Repeating an unkeyed affirmative input at a
new interrupt is a fresh input from the caller. Authorized edits retain
the earlier proposal and approval payload in checkpoint history.
These conditions characterize routing and explicit-edit semantics for
the tested interfaces. They do not exercise same-identifier pending-state
substitution within a shared interrupted thread, Agent Server
authorization policies, or other persistence backends. Matching an
application grant to a candidate during queued handoffs or reuse is a
separate design question.
\texttt{RAW-EVIDENCE.json} preserves the framework identifiers,
requests and checkpoint history; \texttt{OBSERVATIONS.json} records
the relationships compared across runs.

\textbf{X2: binding and execution checks.}
Nine constructed fixtures cross three main binding comparators and one
raw-content diagnostic with two common execution modes: 54 main routes
and 18 diagnostic routes per run. The comparators share a request-value guard;
the content projection extends the retained context with decision-relevant
fields. Amount normalization follows the experiment's declared two-decimal
rule. These selected comparators do not implement the manuscript's full
$G$, nor do they replace its canonical-JSON equality.
The unchecked executor reads the external document at execution; the
verified executor compares its bytes to the reviewed digest and consumes
those same bytes. Both modes are applied to every comparator.

For context, normalized content, instance and raw-content binding,
respectively, semantic-violation counts are $4,1,1,1$ with unchecked
execution and $3,0,0,0$ with verified execution (nine fixtures per cell).
Corresponding SQLite-write counts are $8,5,2,4$ and $7,4,1,3$, totaling
34 effects per run. The three remaining context violations concern fields
omitted from its retained comparison. The changed-document case is
stopped by the shared executor check in every arm. Instance binding also
rejects three same-semantics submissions with new identifiers, recording
the cost of the selected attribution rule separately from semantic
violations. The scorer reads the saved approved-effect fixtures and
SQLite in read-only mode without importing the gate or its projection.

\textbf{Saved runs and verification.}
The archive's \texttt{v2/runs/final-1/} and
\texttt{v2/runs/final-2/} contain runner metadata, checks, raw evidence,
scores and SQLite databases. Their semantic JSON/CSV and actual SQL rows
agree. Randomized framework identifiers are retained but not compared
byte for byte. The original \texttt{VERIFICATION.json} also records
historical source-preservation checks at experiment completion; those
checks describe that time, before the manuscript was revised.
The archive's \texttt{verify\_saved.py} verifies the preserved artifact
hashes, compares the two saved executions and independently rescores
the X2 databases without asserting that the later manuscript remains
unchanged.

\textbf{Reproduction entry point.}
After extracting the archive, \texttt{REPRODUCE.md} gives commands for
read-only validation and optional fresh execution. The saved evidence can
be checked with \texttt{python verify\_saved.py}; fresh execution uses
the pinned requirements and \texttt{v2/run\_all.py} with new output
directories. New runs do not overwrite the archived results. The
unaltered \texttt{v2/FINAL-MANIFEST.json} pins 111 source and result files
and has SHA-256
\nolinkurl{7e850e40e45f2afeffed1a2ab5eadb62738343043406049f5a35dc09d55cb066}.
These deterministic checks complement the separate paired-agent study;
they measure interface and binding behavior rather than autonomous
recovery or human-review costs.

\section{Archived recovery availability and matched task routes}
\label{sec:a12-availability}

This post-hoc readback adds no trajectories to the 32 delivered G pairs. It asks
whether the chosen reuse and reauthorization routes were associated with
different availability of the reviewed candidate and original grant at the
saved rejection boundary. The script \texttt{qa/recovery\_availability.py}
reads the same two frozen archives as the recovery-cost decomposition and
checks their whole-file SHA-256 hashes before reading any episode.

For each episode, the script resolves the saved delivered-feedback event and
verifies the hashes of its rejection-state and model-input references. It
independently checks thirteen predicates: active unconsumed grant, active
reviewer, candidate digest, context digest, preservation of the reviewed
candidate, predecessor freshness, evidence bytes, evidence locator,
source-supported proposal, retained metadata, matching record/field target,
matching review context, and visible grant/candidate/authorized value.
Here, visible means present in the saved agent input, not assessed human
perception or understanding. The script does not use \texttt{legal\_paths\_at\_rejection} to decide availability and
does not import the admission engine. The rejected attempt preserves the
business head and transition list; request and receipt bookkeeping may change.
The observed subsequent commit is checked against the hashed before/after
snapshots and must add exactly one transition.

All thirteen predicates hold in all 32 episodes. Each visible review response
contains the original candidate, authorization identifier and authorized
value. The first post-feedback action is a commit in the 24 reuse episodes,
or an authorization request followed by a commit in the eight reauthorization
episodes. Thus a visible valid reuse option coexists with every observed
reauthorization choice. This is an availability statement at a recorded
boundary, not a counterfactual replay, an optimality claim or a guarantee
against state changes during a production delay.

\begin{table}[htbp]
\centering\small
\caption{Matched task routes at the 16 G task definitions. Reuse means committing the reviewed candidate under the existing grant; fresh means requesting matching new authorization. All 32 episodes have a visible valid reuse option.}
\begin{tabular}{llll}
\toprule
Task suffix & A collection & A route & Original B route\\
\midrule
001 & Supplementary & Reuse & Reuse\\
002 & Supplementary & Reuse & Reuse\\
003 & Supplementary & Reuse & Reuse\\
004 & Original & Fresh & Reuse\\
005 & Supplementary & Reuse & Reuse\\
006 & Supplementary & Fresh & Reuse\\
007 & Supplementary & Fresh & Reuse\\
008 & Supplementary & Reuse & Reuse\\
009 & Supplementary & Fresh & Reuse\\
010 & Supplementary & Reuse & Reuse\\
011 & Supplementary & Fresh & Reuse\\
012 & Supplementary & Reuse & Reuse\\
013 & Supplementary & Reuse & Reuse\\
014 & Supplementary & Fresh & Reuse\\
015 & Supplementary & Fresh & Reuse\\
016 & Supplementary & Fresh & Reuse\\
\bottomrule
\end{tabular}
\end{table}

The task-wise comparison is descriptive. A's sixteen episodes combine one
original episode with fifteen later episodes selected by the original quota
interruption; B's sixteen are original. The same task definition does not
make provider, collection period, sampled prefix or hidden serving state
exchangeable. No model-effect significance test is inferred from the eight
discordant route pairs. The original three cohort ledgers remain unchanged.

Outputs \texttt{tables/availability/episodes.csv},
\texttt{matched-tasks.csv} and \texttt{summary.json} retain every task route,
archive member, state/input hash and cohort. The command and the two external
archive dependencies are specified in \texttt{REPRODUCE-RECOVERY.md}.

\subsection*{Secondary uncertainty analysis of complete suffix cost}

The script \texttt{qa/scientific\_statistics.py} resamples complete task
pairs separately within each collection (10,000 draws, seed 20261023).
It retains policy pairing and does not pool the A/B observations of a task
as independent tasks. The percentile 95\% intervals are exploratory secondary
analyses of existing records, conditional on a small constructed task set.
They are not deployment-level performance guarantees. Original A has one
pair, for which no interval is estimated.

\begin{table}[htbp]
\centering\small
\caption{Complete G suffix cost, bound minus context. Ranges describe observed pairs; intervals are exploratory paired-bootstrap intervals}
\begin{tabular}{lrrrr}
\toprule
Collection & Pairs & Mean calls & Call range & 95\% interval\\
\midrule
Original B & 16 & 1.00 & $[0,2]$ & $[0.75,1.25]$\\
Supplementary A & 15 & 1.60 & $[-1,4]$ & $[0.87,2.33]$\\
Original A & 1 & 2.00 & $[2,2]$ & Not estimated\\
\bottomrule
\end{tabular}
\end{table}

The harness performs every context-side handoff commit. The bound continuation
must perform a recovery commit after rejection; its full-suffix difference
also includes subsequent confirmation and subtracts context-side agent calls.
Recovery through the first persistent commit uses the selected route's minimum
in all 32 cases, but full-suffix differences do not. Original B's differences
are zero in two pairs, one in twelve and two in two; supplementary A's range
includes a negative difference. No claim of optimal route selection follows.

\section{Exploratory permission-state and feedback comparison}
\label{sec:a13-feedback}

\textbf{Purpose and provenance.}
This study was separately frozen before collection on 14 October 2026
(freeze time 13 October, 17:09:12 UTC) and is incorporated into the present
revision after reauditing its raw records. It asks whether explicit legal-path
labels add value beyond the same visible permission facts, and observes
continuation when permission changes after rejection. It is not a new
collection for this revision. The four task instances, FORMAL-G-003 through
FORMAL-G-006, are the next IDs not used in two preceding feedback minipilots.
They reuse the original B review checkpoints from the existing G task family;
they are not unseen tasks relative to the full research program. Selection
does not use their original suffix outcomes. Earlier minipilot results are
not pooled with this comparison.

\textbf{Design.}
Each task has three administrative states after the common identity-mismatch
rejection: \emph{valid} (old grant and reviewer usable), \emph{revoked}
(old grant revoked, reviewer usable), and \emph{unavailable}
(reviewer unavailable, no legal path). Candidate content, evidence, target
and predecessor remain valid. Two independent arms receive either
\emph{facts}, the historical rejection and current permission facts, or
\emph{facts plus options}, those identical facts plus derived legal-path
labels. Removing the options field leaves the initial visible messages
byte-identical; initial states also match in all twelve pairs. The deterministic
mapping uses only common visible facts, not hidden goals or a scoring oracle.
It is valid for this fixture, not a general concurrent authorization planner.

The twelve task--state pairs run in a seeded order (20261014), with arms
adjacent and each state balanced two pairs per starting arm. There are 24
hosted trajectories and 24 separate simple-rule trajectories, one realization
per cell. The requested model alias is \texttt{deepseek-v4-flash} through
the same provider interface as B, at temperature zero. The output cap is
393,216 tokens, the request timeout 1,800 seconds, the episode limit 21,600
seconds and the decision limit six; one transient-transport retry is allowed.
The returned alias does not identify an immutable serving revision. These
budgets and the collection period differ from the original B study.

\textbf{Endpoints and complete results.}
The primary comparison is subsequent tool calls and rejected operations,
options minus facts. Feasible completion and unavailable-state stopping
have separate denominators. Unavailable episodes have $U=0$ because the
task is not committed; a correct stop requires unchanged business state,
no new transition and a truthful report, and is not counted as completion.

\begin{table}[htbp]
\centering\small
\caption{Exploratory feedback results. Each column contains twelve hosted episodes across four tasks and three permission states. Counts are totals, not independent-task estimates.}
\begin{tabular}{lrr}
\toprule
Observation & Facts & Facts plus options\\
\midrule
Feasible task completion & 8/8 & 8/8\\
Unavailable-state correct stops & 4/4 & 4/4\\
Integrity $I=1$ & 12/12 & 12/12\\
Tool calls & 16 & 16\\
Commits / authorization requests & 8/4 & 8/4\\
Rejected continuation operations & 0 & 0\\
Model requests / responses & 28/28 & 28/28\\
Elapsed seconds, total & 371.048 & 283.751\\
Output tokens & 66,774 & 49,293\\
Provider input tokens & 9,747 & 9,376\\
Cache-read input tokens & 70,400 & 72,064\\
Cache-creation input tokens & 0 & 0\\
\bottomrule
\end{tabular}
\end{table}

The twelve paired tool-call differences are $-1$ in three, zero in six and
$+1$ in three; task-level sums are 0, 0, $-1$ and $+1$. Facts uses three
state queries and one receipt query; options uses four receipt queries.
Grouped by permission state, the four paired call differences sum to $+2$
for valid grants, $+1$ for revoked grants and $-3$ for unavailable review.
All three negative differences omit a state query when no legal route remains;
all three positive differences add a receipt query in feasible episodes.
These are descriptive contrasts across four reused tasks per state, not
estimates of a feedback-by-state interaction. Table~\ref{tab:feedback-strata}
retains the corresponding secondary time and output-token totals.

\begin{table}[htbp]
\centering\small
\caption{Secondary feedback outcomes by permission state. Each row contains four pairs; F is facts and L is facts plus labels}
\label{tab:feedback-strata}
\begin{tabular}{lrrrrrr}
\toprule
 & \multicolumn{2}{c}{Calls} & \multicolumn{2}{c}{Seconds} & \multicolumn{2}{c}{Output tokens}\\
State & F & L & F & L & F & L\\
\midrule
Valid & 5 & 7 & 90.407 & 97.813 & 15,773 & 16,863\\
Revoked & 8 & 9 & 155.578 & 142.954 & 27,799 & 24,963\\
Unavailable & 3 & 0 & 125.063 & 42.984 & 23,202 & 7,467\\
\bottomrule
\end{tabular}
\end{table}
These legal reads are confirmation behavior, not rejected or invalid actions.
The simple-rule controller produces identical action sequences across the
two feedback variants, completes all sixteen feasible rule episodes and
correctly stops in all eight unavailable rule episodes. It always confirms
the final state and therefore uses 24 calls per feedback variant. It is not
an optimal controller; the smaller model call total does not establish an
LLM-specific advantage.

All 56 hosted responses terminate naturally. Four exceed 4,096 output
tokens; the largest has 6,522. No response is truncated, and no retry,
transport failure, timeout, quota failure or unfinished episode occurs.
Elapsed time includes feedback construction, recording and continuation;
it excludes preparation, offline audit and human review. Time decreases in
eight of twelve pairs (median paired difference $-7.219$ seconds), and
output tokens decrease in nine. These secondary descriptive directions do
not establish a stable overall cost advantage when the primary tool-call
outcome has no net improvement.

Options are computed without a model call. Across twelve hosted notices,
the options computation takes 37,600 ns in total, and adds 31, 61 or 78
UTF-8 bytes per notice. This in-memory fixture does not measure production
database access, concurrent consistency or rule maintenance. Provider usage
fields remain separate and are not converted to unverified monetary cost.

\textbf{Interpretation and reproduction.}
The comparison supports permission-sensitive continuation under the tested
instructions and exposes no net tool-call benefit from path labels.
It is not an equivalence test, a prevalence estimate, evidence about real
review queues or a cross-model comparison. The four tasks share one family;
repeated states do not create twelve independent task samples.

The supplied \texttt{Feedback-sensitivity-reproduction.zip} retains the
frozen study directory, its pinned source dependencies and a read-only
entry point. \texttt{qa/feedback\_sensitivity.py} verifies eighty frozen
files, reconstructs all 48 hosted/rule episodes from saved inputs,
responses and state transitions, and recomputes the descriptive analysis.
All 48 episode audits pass. The 1,369 files in the study directory remain
unchanged by the readback. Source hashes, cell results and all paired cost
differences are retained in \texttt{tables/feedback-sensitivity/}; the archive's
\texttt{REPRODUCE.md} specifies a fresh output directory and no provider calls.

The source-package script \texttt{qa/feedback\_strata.py} groups the saved
model pairs by permission state and writes \texttt{strata.json}. This grouping
adds no trajectories and leaves the frozen experimental archive unchanged.

\section{Native interrupt and relational contract composition}
\label{sec:a14-composition}

\textbf{Question and design.}
This deterministic experiment connects the native interruption mechanism and
the relational contract in one scripted execution. It tests whether an
application adapter can preserve the reviewed-candidate relation through
approval, handoff, rejection and continuation. It uses LangGraph 1.0.8,
langgraph-checkpoint 4.0.0, langchain-core 1.2.13 and SQLAlchemy 2.0.23 with
Python 3.11. The existing \texttt{stores.py}, its source manifest and nineteen
contract source files are copied without edits. Native framework code is also
unmodified. The protocol, runner, independent scorer, scorer controls and
dependencies are frozen as 27 files before two final executions, each with
fresh SQLite databases. The development round remains separate.

\textbf{Adapter responsibilities.}
An author-written graph first pauses over an already registered immutable
candidate proposing \texttt{machine proposal} for the address field. The
scripted reviewer can deny or authorize \texttt{reviewed correction}; only
an affirmative response to the current native interrupt creates the
application grant handle. A second interrupt is an instrumentation barrier
for injecting the controlled handoff, not another approval request.
The graph invokes the unchanged relational gate with a complete two-field
successor. After rejection, a recovery interrupt receives a scripted choice.
Reuse submits the old candidate/grant; fresh authorization visits a new
review interrupt before creating a grant for the replacement. The unchanged
batch field supplies a source-preservation check. X3 assumes the trusted,
content-immutable candidate storage required in main-text Section 3.2. It
tests the resulting identity relations at commit; it does not independently
probe same-ID content mutation or object replacement between check and use.

The graph, payload interpretation, grant-handle registry, interventions and
recording are adapter code. InMemorySaver and pre-commit grant handles are
in process. The database atomically persists the admitted successor and its
authorization/transition relations. This design does not provide or test
restart-safe approval storage, authenticated review or a distributed
transaction between checkpoint and business storage.

\begin{table}[htbp]
\centering\small
\caption{All ten composition conditions, with identical semantic observations in both final executions. Counts are per execution. ``No write'' refers to business versions and transactional relations, not candidate registration.}
\begin{tabularx}{\textwidth}{@{}>{\raggedright\arraybackslash}p{0.29\textwidth}Xr@{}}
\toprule
Condition & Observed behavior & Address commits\\
\midrule
Corrected approval & Current approval commits corrected value with reviewed identity. & 1\\
Denial & No grant handle or successor write. & 0\\
Identity substitution; reuse & First attempt rejects; old reviewed candidate/grant commits. & 1\\
Identity substitution; fresh review & First attempt rejects; old interrupt ID cannot approve replacement; current ID permits it. & 1\\
Wrong executed value & First attempt rejects; original correction commits on reuse. & 1\\
Stale predecessor & Independent batch update persists; stale address attempt rejects. & 0\\
Rollback after compare-and-swap & No partial transactional rows; retry commits once. & 1\\
Rollback before commit & No partial transactional rows; retry commits once. & 1\\
Duplicate native resume & Completed execution has no additional write. & 1\\
Consumed application-grant replay & Direct replay has no additional write. & 1\\
\bottomrule
\end{tabularx}
\end{table}

\textbf{Independent readback.}
Every graph invocation and gate attempt retains a SQLite copy and the native
pending interrupt identifiers, payloads and scripted input. The scorer
imports neither the gate nor LangGraph. It opens fresh read-only SQLite
connections and checks version counts, corrected values, the retained batch
source, and joins among candidate certificates, authorization bindings,
human-decision rows and fact transitions. It also compares the complete
record-version, transition, decision and authorization-binding tables before
and after every rejected or rolled-back attempt. Gate-return status is not
used to determine whether a business effect occurred. Candidate registration
is permitted before rejection and is kept separate from transactional
business effects.

Both executions satisfy all ten conditions. Each contains fourteen gate
attempts: eight admitted corrections and six rejected or rolled-back attempts
with no partial transactional rows. The stale case contributes one separate
independent batch update and no address correction. Repeating a completed
native resume or explicitly replaying a consumed application grant adds no
business version. In the fresh-review case, the new native interrupt ID
differs from the original; submitting the original ID leaves the current
review pending and creates no new grant or write. Submitting the current ID
permits the replacement's corrected transition with a matching authorization.

Four negative controls mutate database copies: an incorrect final value,
a deleted transition, an incorrect authorized value and an extra version.
The independent scorer detects all four. These checks exercise the oracle;
they are not additional task successes or security-attack trials.

\textbf{Interpretation and artifacts.}
The evidence establishes compatibility of native interrupt routing and the
existing relational checks for the specified adapter and scripted paths.
It does not establish autonomous recovery, production performance, concurrent
replay resistance, restart durability, human review validity or protection
against hostile mutation of pending graph state. Two executions test
deterministic repeatability on the same machine, not statistical robustness.
The agent trajectories reported elsewhere remain a separate evidence layer.

\texttt{Native-contract-reproduction.zip} contains \texttt{PROTOCOL.md},
\texttt{FREEZE.json}, \texttt{run.py}, \texttt{score.py},
\texttt{test\_scorer.py}, the unchanged backend, requirements, both final
run directories and all SQLite snapshots. \texttt{REPRODUCE.md} gives
read-only rescoring and fresh-directory rerun commands. No provider access
or model calls are needed. Development records are retained separately in
the local experiment directory and are not pooled with final observations.

\section{Obligations and existing foundations}
The table retains the detailed mechanism-to-obligation mapping. Bibliographic
entries for the named mechanisms are provided in the main manuscript's related
work. The mapping identifies reused foundations and the role of each obligation
in this controlled comparison; it does not assert that existing systems cannot
express the same relations.

\begin{table}[tbp]
\caption{Existing foundations and their role in the corrected-execution comparison. P1--P6 are defined in main-text Section 3.2.}\label{tab:obligation-positioning}
\centering\small
\begin{tabularx}{\textwidth}{@{}>{\raggedright\arraybackslash}p{0.12\textwidth}>{\raggedright\arraybackslash}p{0.37\textwidth}>{\raggedright\arraybackslash}X@{}}
\toprule
Obligation & Existing basis & Role in this study \\
\midrule
P1 & Continuity Kernel binds proposal and sealed candidate; ReSA tracks the consumed object. & Select context or reviewed-instance agreement as the binding condition while holding correction and shared checks fixed; observe continuation. \\
P2 & ReSA requires preserved or atomically validated consumed versions. & Check internal record, field and evidence consistency against the submitted immutable candidate. Cross-candidate equality belongs to P1; external referents need a separate consumption check. \\
P3 & Guarded Commits binds authority and pre-edit, delta and post-edit hashes. & Check principal authority and the explicit authorized value, retaining the machine proposal separately from the correction and committed value. \\
P4 & Optimistic concurrency control and Continuity Kernel validate predecessors. & Keep freshness shared across policies, separate from instance mismatch. \\
P5 & Transactions and Continuity Kernel install an atomic accepted unit. & Install the complete corrected successor, sources, transitions and grant consumption together. \\
P6 & Provenance and curated-database copying retain derivation and source relations. & Link changed fields to authorized transitions and copy unchanged fields' exact prior sources; check these relations in persistence. \\
\bottomrule
\end{tabularx}
\end{table}

P1 is policy-relative: context binding requires $K_{\mathrm{eq}}$, and bound
binding adds $I_{\mathrm{eq}}$. P2 is context/evidence integrity and P3 is
principal/authorized-value integrity in both policies. The source-level Alloy
names use a historical numbering scheme; Section~\ref{app:formal-mapping}
gives their precise correspondence and verification limits.

The comparison's admission-set inclusions follow from the declared predicates.
The store checks assess implementation of the shared obligations; the agent
trajectories assess choices and costs after admission or rejection. The
observed reuse, renewed authorization and confirmation actions cannot be
derived from the binding predicates alone.

\section{Dynamic recovery after an initially available route}
\label{sec:dynamic-recovery}
This separately frozen exploratory pilot distinguishes route availability at
feedback, authorization at execution, and the eventual business effect. It reuses
the bound Host and independent oracle without modification, and two preselected
checkpoint definitions (FORMAL-G-003 and FORMAL-G-004) from the earlier feedback
study. Each initially has a valid reviewed candidate, active grant and available
reuse, reauthorization and fresh-proposal routes. The initial visible state is
identical across the four conditions of a task. Eight hosted continuations, eight
adaptive-rule continuations and eight cached-route controls were executed on
23 October 2026 after freezing the protocol, inputs, code, allocation and limits.
They are not added to the original 32 G episodes.

Immediately after the controller chooses its first continuation commit, but before
the Host executes it, the harness either leaves state unchanged, revokes the grant
named by that action, advances the predecessor, or revokes active grants and
disables the reviewer. Selected-grant revocation is an adaptive, controlled
intervention, not a sample of naturally occurring administrative changes.
After the triggered return, all three controllers receive the same current
administrative facts. In the predecessor condition, the existing Host advances the
version while preserving the current reference-note value; a fresh candidate and
authorization can therefore restore a legitimate path. The unavailable-reviewer
condition has no legal continuation route.

All eight hosted first commits were admissible before intervention; only the two
static commits remained admissible at execution. The two revocations and two
unavailable-reviewer cases returned \texttt{AUTHORIZATION\_INVALID}; the two
predecessor changes returned \texttt{STALE\_VERSION}. These six deterministic
rejections are host safety behavior, not autonomous-recovery successes.

\begin{table}[htbp]
\centering\small
\caption{Dynamic pilot: independently read-back business recovery. Each controller
has two trajectories per condition; the controllers are not pooled.}
\begin{tabular}{lrrr}
\toprule
Condition & Hosted & Adaptive rule & Cached route\\
\midrule
Static & 2/2 & 2/2 & 2/2\\
Selected grant revoked & 2/2 & 2/2 & 0/2\\
Predecessor advanced & 2/2 & 2/2 & 0/2\\
No legal route & 0/2 & 0/2 & 0/2\\
\bottomrule
\end{tabular}
\end{table}

The hosted and adaptive-rule controllers reauthorized after revocation and rebuilt
the candidate and authorization after a predecessor advance. All three controllers
correctly stopped in both no-legal-route cases, with unchanged target values and
accurate observed reports. The cached-route controller also stopped after the
other dynamic rejections, leaving four feasible corrections incomplete. A fresh
SQLite connection independently read the persisted host state: all 24 trajectories
had integrity one, and all eight hosted final reports matched that state. Negative
controls rejected a forged success report and unrelated-field damage, and returned
unknown when independent readback was removed. All 24 saved trajectories were
reproduced exactly by offline replay.

\paragraph{A preserved measurement limitation.}
The frozen scorer additionally defines \texttt{task\_completion} by requiring a
successful \texttt{get\_state} at the final version, alongside business recovery
and an accurate report. This \emph{state-query-confirmed completion} measure is
narrower than the original task's allowance for receipt or state evidence. It is
0/6 in feasible hosted continuations: five instead queried a successful receipt,
and one used the successful commit response. All six business corrections are
independently verified and accurately reported. The frozen zeros are retained,
not interpreted as six failed business tasks. Adaptive rules score 6/6 and cached
routes 2/6 on this narrower measure. The pilot thus also exposes a measurement
dependence on confirmation-tool choice; subsequent studies should measure
completion from independent authoritative state and truthful reporting, with the
query choice reported separately.

The hosted collection used 35 requests and 27 continuation tool calls; adaptive
rules used 28 calls and cached routes 16. All responses returned the alias
\texttt{deepseek-v4-flash} and ended naturally, with no truncation, invalid output,
transport retry or budget exhaustion. Eight decisions, a 32,768-token output cap,
120 seconds per request and 900 seconds per episode were fixed before collection.
The raw provider-usage fields and local elapsed times are retained in
\texttt{Dynamic-recovery-reproduction.zip}; they do not estimate production latency.
Two reused task definitions, one realization per cell, controlled interventions
and a mutable service alias do not support statistical generalization, a
feedback-format comparison or LLM superiority. Matching adaptive-rule outcomes
show that the fixture is solvable without an LLM. The evidence establishes a
concrete temporal separation between an available route and a subsequently
authorized commit, and between gate rejection and actual recovery.

\section{Reviewed-argument delivery across a process boundary}
\label{sec:a17-correction-consumption}

This section reports a separately frozen check of the trusted-adapter premise used in
the native-default and native-to-store probes. Those probes assume that a complete
reviewed argument object, supplied per pending tool call, is the object the native
runtime executes. The check asks whether three pinned SDKs deliver each reviewed
object to its intended pending call, and leave the other pending calls unchanged,
when the run is paused for review and then continued, including after a clean process
exit and a new interpreter resumes the persisted state.

The question is deliberately narrow. The application merges the reviewed patch into a
complete five-field argument object and supplies that object for every pending call,
including calls the reviewer did not change. The study therefore measures
transmission fidelity of the reviewed object, not partial-patch inference or
field-level preservation by the SDK. The deterministic model emits the declared calls
and then stops; no language model is called, and no provider request is made. A local
tool body writes its actual runtime call identifier and actual arguments to SQLite,
and a separate SDK-free scorer reads those rows in read-only mode.

The matrix crosses three frameworks, two lifecycles, seven call scenarios, six review
edits and two payload families, giving 504 constructed cells. Lifecycles are an
in-process continuation and an exit after a native pause whose persisted state is
resumed by a newly launched interpreter. Scenarios are one call; two calls with
distinct original arguments; and two calls with equal original arguments, each edited
in the first, second or both positions. Edits are a text replacement, integer zero,
boolean false, empty string, a three-field change and a same-value confirmation.
Payload families are ASCII and text with Unicode, quotes and newlines. The protocol,
case matrix, adapters, scorer and SDK sources were hashed before collection; a
separately recorded twelve-cell development preflight is excluded from these
denominators, and 23 runtime and oracle controls were exercised before the freeze.

\begin{table}[htbp]
\centering\small
\caption{Reviewed-argument delivery: cells whose persisted effect multiset equals the
reviewed expectations exactly. Each lifecycle exercises one native continuation route
per framework; the cells are constructed and are not independent samples.}
\label{tab:a17-correction-consumption}
\begin{tabular}{llrrr}
\toprule
Framework & Lifecycle & Cells & Exact match & Mismatch\\
\midrule
Agno 3.1.10 & in-process continuation & 84 & 24 & 60\\
Agno 3.1.10 & exit and fresh-process resume & 84 & 84 & 0\\
LangChain 1.4.7 / LangGraph 1.2.16 & in-process continuation & 84 & 84 & 0\\
LangChain 1.4.7 / LangGraph 1.2.16 & exit and fresh-process resume & 84 & 84 & 0\\
Pydantic AI 2.60.0 & in-process continuation & 84 & 84 & 0\\
Pydantic AI 2.60.0 & exit and fresh-process resume & 84 & 84 & 0\\
\midrule
Total & & 504 & 444 & 60\\
\bottomrule
\end{tabular}
\end{table}

All 60 mismatching cells belong to one path, the Agno in-process continuation, and
they share one mechanism. In every one of them the first pending call executed with
the complete reviewed arguments supplied for the second pending call, while the second
call executed its own; the fields that differ follow the edit, so the pattern is not
restricted to calls with equal original arguments. Of the 72 multi-call cells on that
path, 60 mismatch; the 24 matching cells are the twelve single-call cells and twelve
cells whose two reviewed payloads coincide, where sharing cannot change the delivered
values. The public issue \#8546 already describes shared human-input state across
pending calls in this framework, so the matrix reproduces a known mechanism and
characterizes its current-version scope. It is not reported as a discovered defect,
not confirmed by the maintainers here, and not counted once per cell.

The in-process arm also differs from the fresh arm in its continuation extension point
and in the object it continues from. A 54-cell ablation in the same frozen family
(Agno only) separates the two. Holding the pause, the reviewed objects and the
recorded arguments fixed, the original-run-object route
(\texttt{continue\_run(run\_response=...)})
mismatched in 12 of 18 cells (all twelve multi-call cells; the six single-call cells matched), while the run-identifier route
(\texttt{continue\_run(run\_id, session\_id, requirements)}) matched in all 18 cells in
the same process and in all 18 cells after a clean exit and a new interpreter. The
mechanism therefore follows the continuation route and the state it continues from, not
the process boundary, and the process contrast alone shows no difference. This is a
two-contrast decomposition rather than a single-factor isolation, because the
entrypoint and object identity change together in the first contrast. The matching
paths are equally bounded: LangChain and Pydantic AI use the same documented
continuation route in both lifecycles, so their arms differ mainly in process and
serialization, and one route per lifecycle is exercised.

The primary observable is the exact multiset of call identifier, tool name and
complete typed payload, compared with the reviewed expectations. Field-level
diagnostics locate the same differences and carry no additional detection power; the
two agreed on every cell. No setup mismatch, unsupported path, collection failure,
worker error, native framework error or blocked network attempt occurred. The
constructed cells are not an ecosystem sample: three frameworks were chosen for prior
support, each lifecycle exercises one route, and no prevalence estimate, confidence
interval or significance test applies. What the check supports is conditional and
descriptive: on the documented routes of two SDKs and on the Agno fresh-process route,
a complete reviewed object reached its intended pending call and the unintended call
was unchanged; on the Agno route that continues from the original in-memory run object
with more than one pending call, the reviewed values of one call were executed by
another, while the run-identifier route was clean in the same process.

\paragraph{Same-identifier substitution.}
A third frozen check (16 cells, Pydantic AI 2.60.0) applies a controlled injection
between the review and the continuation: the state the continuation executes from is
replaced while the native call identifier stays unchanged. When the human decision only
approves the pending call, the continuation
executed the substituted arguments in all four cells that changed the supplied message
history, in both lifecycles; substituting the deferred request object alone changed
nothing (two of two executed the reviewed object). When the decision re-supplies the
reviewed object through the native override argument, the reviewed object was executed
in all eight override cells: six with substitution and two unchanged controls. These three contrasts are conditional on the tested path: matching native call
identifiers did not keep the executed arguments equal to the reviewed ones, while
supplying the complete reviewed object through the native override argument left the
executed arguments unchanged under the same substitution. The framework's
documentation describes client-supplied message history as input outside the boundary
the framework enforces, and does not present native approval as a substitute for
server-side authorization
(\url{https://pydantic.dev/docs/ai/core-concepts/message-history/}, accessed
2026-10-24). The result is therefore an empirical statement about an
application-maintained responsibility boundary, not a framework defect; post-approval
substitution itself is reported in prior work on approval hijacking.

\paragraph{Artifact locations.}
The three supplied archives contain frozen protocols, code, cases, source seals,
raw per-cell databases, scores and reports. Their original freeze SHA-256 values are:
\nolinkurl{Correction-consumption-reproduction.zip},
\nolinkurl{740ccc4fee766849dfaa9fec529c79ba36372cd521138f57672aa586ec89cf8b};
\nolinkurl{Agno-entrypoint-reproduction.zip},
\nolinkurl{f17ade246c0e92371874b9fc6373461f2f42ca1ff85ec9e0d2c8d95dd8076166}; and
\nolinkurl{Reviewed-instance-substitution-reproduction.zip},
\nolinkurl{b5871cf623af9cb81fa5c17ac034c913345c8b5d6ed0836e29ed14da498d7a66}.
In each extracted study directory, \texttt{python verify\_offline.py} checks package,
source and raw-record hashes and recomputes effects from the case definitions using
only the standard library. It needs neither SDK imports nor the recorded installation
paths. Recollection instead requires pinned SDK environments and configured collectors;
their concurrency ladder and no-overwrite/no-replay guards remain unchanged.
\texttt{REPRODUCE-V4.md} documents these separate workflows and the history intervention.
Neither workflow requires model calls.


\section*{Artifact sources and reproduction}
The following repository-relative directories contain the complete protocols, raw or linked records, independent scoring, and reproduction entry points. \texttt{SOURCE-MANIFEST.json} supplied with the manuscript package pins the report versions and archive hashes used for Sections A1--A8 and the current manuscript artifacts. Sections A11 and A14 additionally identify the self-contained X1/X2 and native-to-store archives and their source manifests. Online Resource 1 remains the reference for the original study evidence.
\begin{itemize}
\item Output audit: \nolinkurl{research/kais-revision-output-audit-2026-10-14/RESULTS.md}.
\item Budget sensitivity: \nolinkurl{research/kais-budget-sensitivity-2026-10-14/RESULTS.md}; offline analysis uses \texttt{analysis.py}.
\item Commit uncertainty: \nolinkurl{research/kais-commit-uncertainty-2026-10-14/RESULTS.md}; \texttt{run\_probe.py} with a fresh output directory.
\item CK mapping: \nolinkurl{research/kais-native-baseline-probe-2026-10-14/RESULTS.md}; \texttt{run\_probe.py} with a fresh output file.
\item OpenTofu integration: \nolinkurl{research/kais-approval-workflow-2026-10-14/RESULTS.md}; \texttt{run\_workflow.py} with a fresh output directory and pinned protocol.
\item Historical chain audit: \nolinkurl{research/authorization-granularity-phase1-2026-10-06/}; controlled-case records: the pinned DKE-supplement archive cited in the main text.
\item N exposure recount: \nolinkurl{revisions/2026-10-18-kais-recalibrated/qa/n_exposure.py} and \nolinkurl{revisions/2026-10-18-kais-recalibrated/qa/n-exposure.csv}; source archives are named with hashes in Section A7.
\item ToolGate development and public-reference comparisons: \nolinkurl{research/kais-toolgate-feasibility-2026-10-20/} and \nolinkurl{research/kais-toolgate-public-reference-2026-10-20/}; sources, frozen runs and the retained diagnostic are mapped in Section A8.
\item Native-approval and binding-scope comparisons: \nolinkurl{research/ase-x1-x2-2026-10-21/v2/}; the supplied \texttt{X1-X2-reproduction.zip} retains the corrected sources and both final runs, with instructions in \texttt{REPRODUCE.md}.
\item Reviewed-argument delivery: \nolinkurl{research/correction-consumption-study-2026-10-23/}; A17 identifies its archive and two companion checks. Each supplies raw effects and portable offline verification instructions in \texttt{REPRODUCE-V4.md}.
\end{itemize}
Official workflow sources, archived with hashes in the OpenTofu comparison directory, are \url{https://developer.hashicorp.com/terraform/tutorials/automation/automate-terraform} and \url{https://opentofu.org/docs/cli/commands/apply/}. The CK paper is already cited in the main manuscript; its pinned model is \url{https://github.com/openkedge/pci/blob/74841530988f1324bd927da15519486509372917/src/transactional_continuity_kernel/bounded_model.py}.

The reported rebuilds used new directories on the same machine. Independent-environment validation is retained separately: a clean-environment same-machine rebuild (\nolinkurl{research/kais-clean-reproduction-2026-10-14/}) and a cross-machine Linux rebuild with Python 3.11.13 (\nolinkurl{research/kais-server-reproduction-2026-10-14/}). The post-submission studies in this addendum were not themselves rerun cross-machine.
\end{document}
